\documentclass[a4paper]{article}

%%%% Packages %%%%
\usepackage{amsfonts}
\usepackage{amsmath}
\usepackage{amssymb}
\usepackage{amsthm}
\usepackage{bbm}
\usepackage{booktabs}
\usepackage{calc}
\usepackage{caption}
\usepackage{centernot}
\usepackage{circuitikz}
\usepackage{dsfont}
\usepackage{esint}
\usepackage{esvect}
\usepackage{float}
\usepackage{fullpage}
\usepackage{graphicx}
\usepackage{hhline}
\usepackage{hyperref}
\usepackage{ifthen}
\usepackage{lipsum}
\usepackage{marvosym}
\usepackage{mathrsfs}
\usepackage{mathtools}
\usepackage{microtype}
\usepackage{multicol}
\usepackage{parskip}
\usepackage{siunitx}
\usepackage{svg}
\usepackage{tikz}
\usepackage{tikz-cd}
\usepackage{titlesec}
\usepackage{titling}
\usepackage{wrapfig}
\usepackage{xcolor}
\usepackage[inline]{enumitem}
\usepackage[makeroom]{cancel}
\usepackage[most]{tcolorbox}
\usetikzlibrary{angles}
\usetikzlibrary{calc}
\usetikzlibrary{patterns}
\usetikzlibrary{decorations.pathmorphing}
\usetikzlibrary{decorations.pathreplacing}
\usetikzlibrary{decorations.markings}

\setsvg{inkscapeexe="C:/Program Files/Inkscape/bin/inkscape.exe"}

\hypersetup{
	pdfborder={0 0 0}
}
\makeatletter
\renewcommand{\@makefnmark}{\hbox{\textsuperscript{\color{red}\normalfont\@thefnmark}}}
\makeatother
%%%% Packages %%%%
%%%% Custom Commands %%%%

\DeclareMathOperator{\Div}{div}

% Define a problem counter that resets with the section
\newcounter{problem}

% Redefine the problem numbering to include section and problem numbers
\renewcommand{\theproblem}{\arabic{problem}}

% Define the Problem command
\newcommand{\Problem}{\refstepcounter{problem}\fbox{\textbf{Problem \theproblem:}} }

% Solution Definition
\newcommand{\Solution}[1]{\textbf{Solution\ifthenelse{\equal{#1}{}}{}{ #1}:}}

% This forces LaTeX to reserve vertical space as if every line had a fraction - use in \begin{aligned} .. \end{aligned}
\newcommand{\rowstrut}{\vphantom{\dfrac{1}{2}}}
%%%% Custom Commands %%%%
\input{patterns_L1_17}
\tikzset{
	line/.style={
		line width=1.5pt,
	}
}

\tikzset{
	point/.style={
		circle,
		fill=black,
		inner sep=1.5pt
	}
}

\tikzset{
	spring/.style={
		decorate,
		decoration={
			coil,
			aspect=0.6,
			segment length=7pt,
			amplitude=6pt,
			pre length=5pt,
			post length=5pt
		},
		line width=1.5pt
	}
}

\tikzset{
	bracket/.style={
		decorate,
		decoration={
			brace,
			amplitude=7pt
		},
		line width=1.5pt
	}
}

%%% THIS IS A HELPER GRID FOR TIKZ PICTURES %%%

%\begin{center}
%	\begin{tikzpicture}
%		
%		% Grid
%		\draw[help lines, step=1] (-2,-2) grid (12,12);
%
%		% Axes
%		\draw[->, line width=1] (-2,0) -- (12.5,0) node[right] {$x$};
%		\draw[->, line width=1] (0,-2) -- (0,12.5) node[above] {$y$};
%		
%		% Axis labels
%		\foreach \x in {-2,-1,1,2,...,12}
%		\node[below] at (\x,0) {\small $\x$};
%		
%		\foreach \y in {-2,-1,1,2,...,12}
%		\node[left] at (0,\y) {\small $\y$};
%		
%		% Your drawing
%		\draw[spring] (0,0) -- (10,10);
%	\end{tikzpicture}
%\end{center}



%%% THIS HELPS CHECK HEIGHT/WIDTHS OF PICTURES %%%

%\newsavebox{\mytikz}
%\savebox{\mytikz}{%
%	\begin{tikzpicture}[scale=0.5]
%		...
%	\end{tikzpicture}
%}
%\typeout{HEIGHT = \the\ht\mytikz}
%\typeout{DEPTH  = \the\dp\mytikz}
%\typeout{WIDTH  = \the\wd\mytikz}
%height = \the\ht\mytikz
%\\
%depth = \the\dp\mytikz
%\\
%width = \the\wd\mytikz
%\\
%lines = \the\linewidth

\usepackage[english]{babel}

\usepackage[
left=2.5cm,
right=2.5cm,
top=3cm,
bottom=2.5cm
]{geometry}

\setlength{\droptitle}{-7em}
\title{\Huge Physics Olympiad L1-17}
\author{BigChunguz24 - Last Edit 23/09/2026}
\date{}

\begin{document}
	\maketitle
	
	\vspace{-1.25em}
	\Problem
	% Национална Олимпиада По Физика 2025 - Тема 10 Клас - Задача 3, Част 2
	\textbf{\underline{Spring-Pulley System}}
	\\
	Three massless springs with spring constants $k$, $2k$ and $3k$ are connected by inextensible strings to a horizontal ceiling and a massless pulley as shown. At a certain moment, under the action of an external force $\vv{F}$ that increases smoothly with time, the pulley starts moving downward with a constant velocity $\vv{v}$.
	\begin{description}
		\item[$\bullet$] Find the velocities of points $A, B, C, D$ and $E$;
		\item[$\bullet$] Find the velocity of the string $CE$ relative to the rim of the pulley;
		\item[$\bullet$] Find the time dependence of the force $\vv{F}(t)$ if it is known that $\vv{F}(0) = \vv{0}$;
	\end{description}
	Now assume that the Spring-Pulley system is at rest at height $y=0$, where the $y$-axis points downwards. At $t=0$, a body of mass $m$ is attached to the pulley and released from rest.
	\begin{description}
		\item[$\bullet$] Find the equation of motion $y(t)$ of the body;
		\item[$\bullet$] Find the period $T$ and angular frequency $\omega$ of the oscillations;
		\item[$\bullet$] Find the maximum depth $d$ that the body reaches during this motion;
	\end{description}
	
	\begin{center}
	\begin{tikzpicture}
		\def\shift{6cm}
		
		% Ceilings
		\draw[line] (1, 12) -- (6, 12);
		\draw[line, xshift=\shift] (1, 12) -- (6, 12);
		
		% Ceiling hatchings
		\foreach \x in {1, 1.25, ..., 6}
		\draw[line] (\x, 12) -- ++(0.25,0.25);
		
		\foreach \x in {1, 1.25, ..., 6}
		\draw[line, xshift=\shift] (\x, 12) -- ++(0.25,0.25);
		
		% Pulleys
		\draw[line width=1.5pt] 
		(3.5, 6) circle[radius=0.75];
		
		\draw[line width=1.5pt, xshift=\shift] 
		(3.5, 6) circle[radius=0.75];
		
		% Upper connector (above springs)
		\draw[line] (2.75, 12) -- (2.75, 11);
		\draw[line] (4.25, 12) -- (4.25, 10.5);
		
		\draw[line, xshift=\shift] (2.75, 12) -- (2.75, 11);
		\draw[line, xshift=\shift] (4.25, 12) -- (4.25, 10.5);
		
		% Lower connectors (below springs)
		\draw[line] (2.75, 7) -- (2.75, 6);
		\draw[line] (4.25, 7.5) -- (4.25, 6);
		\draw[line, xshift=\shift] (2.75, 7) -- (2.75, 6);
		\draw[line, xshift=\shift] (4.25, 7.5) -- (4.25, 6);
		
		% Spring points
		\node[point, label=left:$A$] at (2.75, 11) {};
		\node[point, label=left:$B$] at (2.75, 9) {};
		\node[point, label=left:$C$] at (2.75, 7) {};
		\node[point, label=right:$D$] at (4.25, 10.5) {};
		\node[point, label=right:$E$] at (4.25, 7.5) {};
		
		\node[point, xshift=\shift] at (2.75, 11) {};
		\node[point, xshift=\shift] at (2.75, 9) {};
		\node[point, xshift=\shift] at (2.75, 7) {};
		\node[point, xshift=\shift] at (4.25, 10.5) {};
		\node[point, xshift=\shift] at (4.25, 7.5) {};
		
		% Springs
		\draw[spring] (2.75, 11) -- (2.75, 9);
		\draw[spring] (2.75, 9) -- (2.75, 7);
		\draw[spring] (4.25, 10.5) -- (4.25, 7.5);
		
		\draw[spring, xshift=\shift] (2.75, 11) -- (2.75, 9);
		\draw[spring, xshift=\shift] (2.75, 9) -- (2.75, 7);
		\draw[spring, xshift=\shift] (4.25, 10.5) -- (4.25, 7.5);
		
		% Spring Stiffness
		\node[left] at (2.5, 10) {$k$};
		\node[left] at (2.5, 8) {$2k$};
		\node[right] at (4.5, 9) {$3k$};
		
		\node[left, xshift=\shift] at (2.5, 10) {$k$};
		\node[left, xshift=\shift] at (2.5, 8) {$2k$};
		\node[right, xshift=\shift] at (4.5, 9) {$3k$};
		
		% Pulley centres
		\node[point] at (3.5, 6) {};
		\node[point, xshift=\shift] at (3.5, 6) {};
		
		% Vector
		\draw[line, -{Stealth}] (3.5, 6) -- (3.5, 4.5);
		\node[right] at (3.5, 4.5) {$\vv{v}$};
		
		% Mass
		\draw[line, xshift=\shift] (3.5, 6) -- (3.5, 4.5);
		\draw[line, xshift=\shift] (3, 4) rectangle (4, 4.5);
		\node at (9.5, 4.25) {$m$};
	\end{tikzpicture}
\end{center}
	
	\Problem 
	% Национална Олимпиада По Физика 2013 - Тема 10-12 Клас - Задача 2, Част 1
	\textbf{\underline{Metal Ring in Magnetic Field}}
	\\
	A metal ring of radius $r = 0{,}1\,\mathrm{m}$ and resistance $R = 2{,}0\,\mathrm{\Omega}$ is positioned horizontally in a uniform vertical magnetic field of magnitude $B_0 = 5{,}0\cdot10^{-5}\,\mathrm{T}$. It is then flipped over so that its other side faces upward.
	\begin{description}[leftmargin=!, labelwidth=\widthof{$\bullet$}]
		\item[$\bullet$] What magnitude of electric charge $q$ passes through the cross section of the ring?
		\item[$\bullet$] The ring is flipped uniformly in a time $\tau = 1\,\mathrm{s}$. Find the magnitude of the current as a function of time;
	\end{description}
	Thermal effects and self-inductance can be ignored.

	\Problem
	% Национална Олимпиада По Физика 2013 - Тема 10-12 Клас - Задача 2, Част 2
	\textbf{\underline{Counting Waves}}
	\\
	Between two points on a sound wave, which vibrate in phase, there are $N = 825$ wavelengths. When the temperature of the medium is increased by $\Delta T_0 = 1\,\mathrm{K}$, the speed of sound increases by $\eta = 0{,}2\,\%$. 
	\begin{description}[leftmargin=!, labelwidth=\widthof{$\bullet$}]
		\item[$\bullet$] Find the temperature $T$ of the medium;
		\item[$\bullet$] What temperature increase $\Delta T$ is required for the two points to start vibrating in antiphase?
	\end{description}
	The medium is assumed to be an ideal gas and the frequency of the sound remains constant.

	\Problem
	% 3800 Задач Физике - 4.39
	\textbf{\underline{Resultant of Parallel Forces}}
	\\
	Two parallel and equally directed forces $F_1$ and $F_2$ are applied to a rigid body. Prove that:
	\begin{description}[leftmargin=!, labelwidth=\widthof{$\bullet$}]
		\item[$\bullet$] The resultant force is parallel to the applied forces and directed in the same direction;
		\item[$\bullet$] The magnitude of the resultant force is equal to the sum of the magnitudes of the applied forces;
		\item[$\bullet$] The line of action of the resultant passes through a point that divides the segment between the points of application of the two forces into parts that are inversely proportional to the force magnitudes;
	\end{description}
	
	\Problem
	% 3800 Задач Физике - 11.231 + 11.232 + 11.233
	\textbf{\underline{Electron-Ring Interactions}}
	\\
	We consider the motion of an electron along the axis of a uniformly charged ring with linear charge density $\lambda > 0$. Assuming the mass of the ring to be much larger than that of the electron, find:
	\begin{description}[leftmargin=!, labelwidth=\widthof{$\bullet$}]
		\item[$\bullet$] With what speed will the $e^-$ pass through the centre of the ring if it is attracted from infinity?
		\item[$\bullet$] With what speed will the $e^-$ pass through the centre of the ring if it is attracted from distance $\eta R$?
		\item[$\bullet$] What is the period of small oscillations of the $e^-$ about the centre of the ring along its axis?
		\item[$\bullet$] What is the largest amplitude for which the magnitude of the restoring force differs from that predicted by the small-oscillation approximation by no more than $1\,\%$? Specify your answer as a multiple of $R$.
		\item[$\bullet$] What is the largest amplitude for which the period of oscillation differs from that predicted by the small-oscillation approximation by no more than $1\,\%$? Specify your answer as a multiple of $R$.
	\end{description}
	The electron has mass $m_e$ and charge $-e$. The radius of the ring is $R$ and $\eta$ is a real number ($\eta \in \mathbb{R}^+$). In all situations where relevant, the electron starts its motion from rest. You may use appropriate approximations and/or cutoffs where necessary.
	
	\Problem
	% Irodov 2.148
	\textbf{\underline{Free Expansion}}
	\\
	One of the two thermally insulated vessels interconnected by a tube with valve contains $\nu = 2{,}2$ moles of an ideal gas. The other vessel is evacuated. The valve having been opened, the gas increases its volume $n = 3{,}0$ times. Find the entropy increment of the gas.
	
	\newpage
\vspace*{-5em}
\enlargethispage{\baselineskip}
\newcommand{\tg}{\operatorname{tg}}
\newcommand{\arctg}{\operatorname{arctg}}
\section*{\centering Appendix}
\section*{\centering Useful Mathematical Formulae and Physical Constants}
\emph{\underline{Important}}: Infinite series that are marked with $(*)$ are only valid under the condition that $|x| < 1$.

\begin{minipage}[t]{0.5\textwidth}
	\vspace{0pt}
	\begin{tcolorbox}[title=\textbf{Taylor Expansions}]
		\vspace{-10pt}
		\begin{align*}
			& \sin(x) = x - \frac{x^3}{6} + \frac{x^5}{120} + \mathcal{O}(x^7)
			&&
			\\[1pt]
			& \cos(x) = 1 - \frac{x^2}{2} + \frac{x^4}{24} - \frac{x^6}{720} + \mathcal{O}(x^8)
			&&
			\\[1pt]
			& \tg(x) = x + \frac{x^3}{3} + \frac{2x^5}{15} + \mathcal{O}(x^7)
			&&
			\\[1pt]
			& \arcsin(x) = x + \frac{x^3}{6} + \frac{3x^5}{40} + \mathcal{O}(x^7)
			&&
			\\[1pt]
			& \arctg(x) = x - \frac{x^3}{3} + \frac{x^5}{5} + \mathcal{O}(x^7)
			&&
			\\[1pt]
			& \exp(x) = 1 + x + \frac{x^2}{2} + \mathcal{O}(x^3)
			&&
			\\[1pt]
			& \ln(1+x) = x - \frac{1}{2}x^2 + \mathcal{O}(x^3)
			&& (*)
			\\[1pt]
			& (1-x)^{-1} = 1 + x + x^2 + \mathcal{O}(x^3)
			&& (*)
			\\[1pt]
			& (1+x)^{-1} = 1 - x + x^2 + \mathcal{O}(x^3)
			&& (*)
			\\[1pt]
			& (1+x)^{+\frac{1}{2}} = 1 + \frac{1}{2}x - \frac{1}{8}x^2 + \mathcal{O}(x^3)
			&& (*)
			\\[1pt]
			& (1+x)^{-\frac{1}{2}} = 1 - \frac{1}{2} x + \frac{3}{8} x^2 + \mathcal{O}(x^3)
			&& (*)
			\\[1pt]
			& (1+x)^{+\frac{3}{2}} = 1 + \frac{3}{2}x + \frac{3}{8}x^2 + \mathcal{O}(x^3)
			&& (*)
			\\[1pt]
			& (1+x)^{-\frac{3}{2}} = 1 - \frac{3}{2}x + \frac{15}{8}x^2 + \mathcal{O}(x^3)
			&& (*)
		\end{align*}
		
	\end{tcolorbox}
\end{minipage}
\begin{minipage}[t]{0.5\textwidth}
	\vspace{0pt}
	\begin{tcolorbox}[title=\textbf{Inverse Hyperbolic Functions}]
		\vspace{-10pt}
		\begin{align*}
			& \text{arsinh}(x) = \ln(x + \sqrt{x^2 + 1})
			\\[3.0375pt]
			& \text{arcosh}(x) = \ln(x + \sqrt{x^2 - 1})
			\\[3.0375pt]
			& \text{artgh}(x) = \frac{1}{2} \ln
			\left(
			\frac{1+x}{1-x}
			\right)
			\\[3.0375pt]
			& \text{arctgh}(x) = \frac{1}{2} \ln
			\left(
			\frac{x+1}{x-1}
			\right)
		\end{align*}
		
	\end{tcolorbox}
	\begin{tcolorbox}[title=\textbf{Indefinite Integrals}]
		\vspace{-10pt}
		\begin{align*}
			& \int (1-x^2)^{-\frac{3}{2}} dx = \frac{x}{\sqrt{1-x^2}} + C
			\\[3.0375pt]
			& \int (1+x^2)^{-\frac{3}{2}} dx = \frac{x}{\sqrt{1 + x^2}} + C
			\\[3.0375pt]
			& \int (1-x^2)^{-\frac{1}{2}} dx = \arcsin(x) + C
			\\[3.0375pt]
			& \int (1+x^2)^{-\frac{1}{2}} dx = \text{arsinh}(x) + C
			\\[3.0375pt]
			& \int (1-x^2)^{+\frac{1}{2}} dx = \frac{\arcsin(x) + x \sqrt{1-x^2}}{2} + C
			\\[3.0375pt]
			& \int (1+x^2)^{+\frac{1}{2}} dx = \frac{\text{arsinh}(x) + x \sqrt{1 + x^2}}{2} + C
		\end{align*}
		
	\end{tcolorbox}
\end{minipage}

\begin{minipage}{\textwidth}
	\begin{tcolorbox}[title=\textbf{Useful Physical Constants}]
		\vspace{-5pt}
		\begin{center}
			\begin{tabular}{|l|c|c|c|}
				\hline
				\textbf{Constant} & \textbf{Symbol} & \textbf{Value} & \textbf{Measuring Units} 
				\\
				\hline
				Light speed in vacuum & $c$ & $3{,}000 \times 10^8$ & $\mathrm{m/s}$
				\rule{0pt}{10pt}
				\\[1pt]
				\hline
				Rest mass of electron & $m_e$ & $9{,}109 \times 10^{-31}$ & $\mathrm{kg}$
				\rule{0pt}{10pt}
				\\[1pt]
				\hline
				Rest mass of proton & $m_p$ & $1{,}673 \times 10^{-27}$ & $\mathrm{kg}$
				\rule{0pt}{10pt}
				\\[1pt]
				\hline
				Rest mass of neutron & $m_n$ & $1{,}675 \times 10^{-27}$ & $\mathrm{kg}$
				\rule{0pt}{10pt}
				\\[1pt]
				\hline
				Rest mass of muon & $m_{\mu}$ & $1{,}884 \times 10^{-28}$ & $\mathrm{kg}$
				\rule{0pt}{10pt}
				\\[1pt]
				\hline
				Rest mass of tau-lepton & $m_{\tau}$ & $3{,}168 \times 10^{-27}$ & $\mathrm{kg}$
				\rule{0pt}{10pt}
				\\[1pt]
				\hline
				Elementary charge & $e$ & $1{,}602 \times 10^{-19}$ & $\mathrm{C}$
				\rule{0pt}{10pt}
				\\[1pt]
				\hline
				Planck's constant & $h$ & $6{,}626 \times 10^{-34}$ & $\mathrm{Js}$
				\rule{0pt}{10pt}
				\\[1pt]
				\hline
				Boltzmann's constant & $k$ & $1{,}381 \times 10^{-23}$ & $\mathrm{J/K}$
				\rule{0pt}{10pt}
				\\[1pt]
				\hline
				Gravitational constant & $\gamma$ & $6{,}673 \times 10^{-11}$ & $\mathrm{Nm^2/kg^2}$ 
				\rule{0pt}{10pt}
				\\[1pt]
				\hline
				Dielectric permittivity of vacuum & $\varepsilon_0$ & $8{,}854 \times 10^{-12}$ & $\mathrm{F/m}$
				\rule{0pt}{10pt}
				\\[1pt]
				\hline
				Avogadro's number & $N_A$ & $6{,}022 \times 10^{23}$ & $\mathrm{mol^{-1}}$
				\rule{0pt}{10pt}
				\\[1pt]
				\hline
				Rydberg's constant & $R_{\infty}$ & $1{,}097 \times 10^7$ & $\mathrm{m^{-1}}$
				\rule{0pt}{10pt}
				\\[1pt]
				\hline
				Magnetic permeability of vacuum & $\mu_0$ & $4\pi \times 10^{-7}$ & $\mathrm{H/m}$
				\rule{0pt}{10pt}
				\\[1pt]
				\hline
				Universal gas constant & $R$ & $8{,}314$ & $\mathrm{J / (mol \cdot K)}$
				\rule{0pt}{10pt}
				\\[1pt]
				\hline
			\end{tabular}
		\end{center}
	\end{tcolorbox}
\end{minipage}
%	\newpage
\section*{\centering Solutions}
\Solution{1}
Due to the upper strings being inextensible, points $A$ and $D$ must have zero velocity:
\begin{equation*}
	\boxed{v_A = 0}
	\quad\text{and}\quad
	\boxed{v_D = 0}
	\tag{$1$}
	\label{eq:j17_p1_1}
\end{equation*}

We can now draw the equivalent configuration by applying the rule for springs in series:
\begin{equation*}
	\frac{1}{k_{eff}} = \frac{1}{k_1} + \frac{1}{k_2} = 
	\frac{1}{k} + \frac{1}{2k} = 
	\frac{3}{2k}
	\quad\Rightarrow\quad
	\boxed{k_{eff} = \frac{2k}{3}}
	\tag{$*$}
	\label{eq:j17_p1_mock_1}
\end{equation*}

\begin{center}
	\begin{tikzpicture}
		\def\shift{6cm}
		
		% Ceilings
		\draw[line] (1, 12) -- (6, 12);
		\draw[line, xshift=\shift] (1, 12) -- (6, 12);
		
		% Ceiling hatchings
		\foreach \x in {1, 1.25, ..., 6}
		\draw[line] (\x, 12) -- ++(0.25,0.25);
		
		\foreach \x in {1, 1.25, ..., 6}
		\draw[line, xshift=\shift] (\x, 12) -- ++(0.25,0.25);
		
		% Pulleys
		\draw[line width=1.5pt] 
		(3.5, 6) circle[radius=0.75];
		
		\draw[line width=1.5pt, xshift=\shift] 
		(3.5, 6) circle[radius=0.75];
		
		% Upper connector (above springs)
		\draw[line] (2.75, 12) -- (2.75, 11);
		\draw[line] (4.25, 12) -- (4.25, 10.5);
		
		\draw[line, xshift=\shift] (2.75, 12) -- (2.75, 11);
		\draw[line, xshift=\shift] (4.25, 12) -- (4.25, 10.5);
		
		% Lower connectors (below springs)
		\draw[line] (2.75, 7) -- (2.75, 6);
		\draw[line] (4.25, 7.5) -- (4.25, 6);
		\draw[line, xshift=\shift] (2.75, 7) -- (2.75, 6);
		\draw[line, xshift=\shift] (4.25, 7.5) -- (4.25, 6);
		
		% Spring points
		\node[point, label=left:$A$] at (2.75, 11) {};
		\node[point, label=left:$B$] at (2.75, 9) {};
		\node[point, label=left:$C$] at (2.75, 7) {};
		\node[point, label=right:$D$] at (4.25, 10.5) {};
		\node[point, label=right:$E$] at (4.25, 7.5) {};
		
		\node[point, label=left:$A$, xshift=\shift] at (2.75, 11) {};
		\node[point, label=left:$C$, xshift=\shift] at (2.75, 7) {};
		\node[point, label=right:$D$, xshift=\shift] at (4.25, 10.5) {};
		\node[point, label=right:$E$, xshift=\shift] at (4.25, 7.5) {};
		
		\node[point, xshift=\shift] at (2.75, 11) {};
		\node[point, xshift=\shift] at (2.75, 7) {};
		\node[point, xshift=\shift] at (4.25, 10.5) {};
		\node[point, xshift=\shift] at (4.25, 7.5) {};
		
		% Springs
		\draw[spring] (2.75, 11) -- (2.75, 9);
		\draw[spring] (2.75, 9) -- (2.75, 7);
		\draw[spring] (4.25, 10.5) -- (4.25, 7.5);
		
		\draw[spring, xshift=\shift] (2.75, 11) -- (2.75, 7);
		\draw[spring, xshift=\shift] (4.25, 10.5) -- (4.25, 7.5);
		
		% Spring Stiffness
		\node[left] at (2.5, 10) {$k$};
		\node[left] at (2.5, 8) {$2k$};
		\node[right] at (4.5, 9) {$3k$};
		
		\node[left, xshift=\shift] at (2.5, 9) {$\displaystyle\frac{2k}{3}$};
		\node[right, xshift=\shift] at (4.5, 9) {$3k$};
		
		% Pulley centres
		\node[point] at (3.5, 6) {};
		\node[point, xshift=\shift] at (3.5, 6) {};
		
		% Vector
		\draw[line, -{Stealth}] (3.5, 6) -- (3.5, 4.5);
		\node[right] at (3.5, 4.5) {$\vv{v}$};
		
		\draw[line, -{Stealth}, xshift=\shift] (3.5, 6) -- (3.5, 4.5);
		\node[right, xshift=\shift] at (3.5, 4.5) {$\vv{v}$};
		
		% Velocities
		\draw[line, -{Stealth}] (5.5, 9) -- (7.5, 9);
		
		\draw[line, -{Stealth}] (2.5, 6.5) -- (2.5, 5.5);
		\node[left] at (2.5, 6) {$\vv{v}_C$};
		
		\draw[line, -{Stealth}] (4.5, 5.5) -- (4.5, 6.5);
		\node[right] at (4.5, 6) {$\vv{v}_E$};
		
		\draw[line, -{Stealth}, xshift=\shift] (2.5, 6.5) -- (2.5, 5.5);
		\node[left, xshift=\shift] at (2.5, 6) {$\vv{v}_C$};
		
		\draw[line, -{Stealth}, xshift=\shift] (4.5, 5.5) -- (4.5, 6.5);
		\node[right, xshift=\shift] at (4.5, 6) {$\vv{v}_E$};
		
		
	\end{tikzpicture}
\end{center}

When replacing springs in series by an equivalent spring, the total extension of the original springs is equal to the extension of the equivalent spring, and the endpoint velocities remain unchanged. Therefore, we can work directly with $\vv{v}_C$ and $\vv{v}_E$, which are also the velocities of the string segments connected to the pulley. Relative to the pulley, these velocities must have equal magnitudes and opposite directions:
\begin{equation*}
	v_C - v = -(v_E - v)
	\quad\Rightarrow\quad
	\boxed{v_C + v_E = 2v}
	\tag{$**$}
	\label{eq:j17_p1_mock_2}
\end{equation*}
Since the string is massless and inextensible, its tension is the same on both sides of the pulley. Therefore:
\begin{equation*}
	k_{eff} \Delta x_{eff} = 3k \Delta x_3
	\quad\Rightarrow\quad
	\frac{\Delta x_{eff}}{\Delta x_3}
	=
	\frac{3k}{k_{eff}}
	\quad\Rightarrow\quad
	\frac{v_C - v_A}{v_E - v_D}
	=
	\frac{9}{2}
	\quad\Rightarrow\quad
	\boxed{\frac{v_C}{v_E} = \frac{9}{2}}
	\tag{$***$}
	\label{eq:j17_p1_mock_3}
\end{equation*}
By equations (\ref{eq:j17_p1_mock_2}) and (\ref{eq:j17_p1_mock_3}) we conclude that:
\begin{equation*}
	\boxed{v_C = \frac{18v}{11}}
	\quad\Rightarrow\quad
	\boxed{v_E = \frac{4v}{11}}
	\tag{$2$}
	\label{eq:j17_p1_2}
\end{equation*}
This also gives us the velocity of the string relative to the rim of the pulley from (\ref{eq:j17_p1_mock_2}):
\begin{equation*}
	|v_{rel}| = |v_C - v| = |v_E - v| = \frac{7v}{11}
	\quad\Rightarrow\quad
	\boxed{|v_{rel}| = \frac{7v}{11}}
	\tag{$3$}
	\label{eq:j17_p1_3}
\end{equation*}
which has the direction depicted on the figure: on the LHS downward and on the RHS upward.

To find the velocity of point $B$, observe that:
\begin{equation*}
	k_1 \Delta x_1 = k_{eff} \Delta x_{eff}
	\quad\Rightarrow\quad
	\frac{k_1}{k_{eff}} = \frac{\Delta x_{eff}}{\Delta x_1}
	\quad\Rightarrow\quad
	\frac{k}{2k/3} = \frac{v_C - v_A}{v_B - v_A}
	\quad\Rightarrow\quad
	\boxed{\frac{v_C}{v_B} = \frac{3}{2}}
	\tag{$\star$}
	\label{eq:j17_p1_mock_4}
\end{equation*}
Finally, applying (\ref{eq:j17_p1_mock_4}) to (\ref{eq:j17_p1_2}) we get:
\begin{equation*}
	\boxed{v_B = \frac{12v}{11}}
	\tag{$4$}
	\label{eq:j17_p1_4}
\end{equation*}
\newpage
The only forces acting on the pulley are the force $\vv{F}$ and the two tensions $\vv{T}$ on either side:
\\[-0.8\baselineskip]
\begin{minipage}[t][75pt]{\textwidth}
	\vspace{0pt}
	\begin{minipage}[t][72.86pt]{0.15\textwidth}
	\vspace{0pt}
	\centering
	\begin{tikzpicture}[scale=0.5]
		% Pulley
		\draw[line] (6, 6) circle[radius=1];
		\node[point] at (6, 6) {};
		
		% Downward Force
		\draw[line, -{Stealth}] (6, 6) -- (6, 3);
		\node[right] at (6, 4) {$\vv{F}$};
		
		% Tensions
		\draw[line, -{Stealth}] (5, 6) -- (5, 7.5);
		\draw[line, -{Stealth}] (7, 6) -- (7, 7.5);
		
		\node[left] at (5, 7.5) {$\vv{T}$};
		\node[right] at (7, 7.5) {$\vv{T}$};
	\end{tikzpicture}
\end{minipage}
	\begin{minipage}[t][72.86pt]{0.84\textwidth}
		\vspace{7pt}
		Considering that this is the same tension each spring carries:
		\begin{equation*}
			F = 2T = 2 k_3 \Delta x_3(t) = 6k (v_E - v_D) t
			\quad\Rightarrow\quad
			\boxed{F(t) = \frac{24kvt}{11}}
			\tag{$5$}
			\label{eq:j17_p1_5}
		\end{equation*}
		where we have used the extension of the spring with stiffness $k_3 = 3k$.
	\end{minipage}
\end{minipage}

In the second part of the problem rest is assumed at $y=0$ at the initial moment $t=0$. We start by finding the equilibrium position of the pulley with mass and the effective stiffness of this system.

Expression (\ref{eq:j17_p1_5}) gives us the magnitude of the external force required to keep the pulley moving with constant velocity. At every instant, this force balances the restoring force exerted by the springs. However, the appearance of the velocity $v$ of the pulley in (\ref{eq:j17_p1_5}) does not mean that the restoring force depends on the velocity of the pulley. It appears only because, under the prescribed motion, the displacement is \(x(t)=vt\). Thus, rewriting the force in terms of the displacement gives:
\begin{equation*}
	F(t) = F_{res}(x(t)) \equiv \frac{24kx(t)}{11}
	\quad\Rightarrow\quad
	F_{res}(x) = \frac{24kx}{11}
	\tag{$\star\star$}
	\label{eq:j17_p1_mock_5}
\end{equation*}
So, to avoid confusion\footnote{The confusion being: zero velocity at equilibrium $\rightarrow F(t) = 0$ by equation (\ref{eq:j17_p1_5}) $\rightarrow$ restoring force vanishes at equilibrium.}, we emphasize that these are different forces:
\begin{description}
	\item[$\rightarrow$] $F(t)$ in equation (\ref{eq:j17_p1_5}) describes an external force which keeps the pulley at $\vv{v} = \text{const}$;
	\item[$\rightarrow$] $F_{res}(x)$ in equation (\ref{eq:j17_p1_mock_5}) describes the restoring force the springs provide to a displaced pulley;
\end{description}
To find the effective stiffness for the Spring-Pulley system:
\begin{equation*}
	F_{res}(x) = \kappa_{eff} x = \frac{24kx}{11}
	\quad\Rightarrow\quad
	\boxed{\kappa_{eff} = \frac{24k}{11}}
	\tag{$6a$}
	\label{eq:j17_p1_6a}
\end{equation*}
This directly gives us the angular frequency and the period of the small oscillations for the body:
\begin{equation*}
	\left.
	T = 2 \pi \sqrt{\frac{m}{\kappa_{eff}}}
	\quad\Rightarrow\quad
	\boxed{T = \pi \sqrt{\frac{11m}{6k}}}
	\quad\,\,\,\,
	\right\vert
	\quad
	\omega = \sqrt{\frac{\kappa_{eff}}{m}}
	\quad\Rightarrow\quad
	\boxed{\omega = \sqrt{\frac{24k}{11m}}}
	\tag{$6b$}
	\label{eq:j17_p1_6b}
\end{equation*}
\begin{minipage}[t][158pt]{\textwidth}
	\vspace{0pt}
	\begin{minipage}[t][156.6pt]{0.25\textwidth}
	\vspace{0pt}
	\centering
	\begin{tikzpicture}
		% Ceiling
		\draw[line] (1.5, 12) -- (3.5, 12);
		
		% Ceiling hatchings
		\foreach \x in {1.5, 1.75, ..., 3.5}
		\draw[line] (\x, 12) -- ++(0.25,0.25);
		
		% String - Spring - String
		\draw[line] (2.5, 12) -- (2.5, 11.5);
		\draw[spring] (2.5, 11.5) -- (2.5, 9.5);
		\draw[line] (2.5, 9.5) -- (2.5, 9);
		\node[left] at (2.3, 10.5) {$\kappa_{eff}$};
		
		% Mass
		\draw[line] (2, 8.5) rectangle (3, 9);
		\node at (2.5, 8.75) {$m$};
		
		% Axis - Equilibrium
		\draw[line, -{Stealth}] (1, 12) -- (1, 7);
		\node[left] at (1, 7) {$y$};
		
		\draw[line, dashed] (1, 8.75) -- (2, 8.75);
		\node[left] at (1, 8.75) {$eq$};
		
		\draw[line] (0.8, 9.75) -- (1.2, 9.75);
		\node[left] at (0.8, 9.75) {$0$};
		
		\draw[line] (0.8, 7.75) -- (1.2, 7.75);
		\node[left] at (0.8, 7.75) {$2A$};
		
		\draw[line, Stealth-Stealth] (3.25, 8.25) -- (3.25, 9.25);
	\end{tikzpicture}
\end{minipage}
\begin{minipage}[t][156.6pt]{0.74\textwidth}
	\vspace{2pt}
	To find the equilibrium position, we balance the forces:
	\begin{equation*}
		\kappa_{eff} y_{eq} = mg
		\quad\Rightarrow\quad
		\boxed{y_{eq} = \frac{11mg}{24k}}
		\tag{$6c$}
		\label{eq:j17_p1_6c}
	\end{equation*}
	As the body is released from $y=0$, it will pass through the equilibrium position, reach a maximum depth, then travel back up to its original position due to energy conservation. Therefore, the amplitude of the oscillations is exactly the distance between the initial position and the equilibrium position:
	\begin{equation*}
		\boxed{A = y_{eq} = \frac{11mg}{24k}}
		\quad\text{and}\quad
		\boxed{d = 2A = \frac{11mg}{12k}}
		\tag{$7$}
		\label{eq:j17_p1_7}
	\end{equation*}
\end{minipage}
\end{minipage}
To find the equation of motion, consider that about the equilibrium we have ($y^*$ and $y$ both point down):
\begin{equation*}
	y^*(t) = - A \cos(\omega t)
\end{equation*}
where it is easy to confirm the initial conditions i.e. $y^*(0) = -A$ and $v(0) = 0$, as required. Finally, the coordinate $y$ is just displaced by one amplitude with respect to the coordinate $y^*$:
\begin{equation*}
	y(t) = y^*(t) + A
	\quad\Rightarrow\quad
	\boxed{y(t) = A(1 - \cos(\omega t))}
	\quad\text{or}\quad
	\boxed{y(t) = \frac{11mg}{24k}
		\left(
		1 - \cos\left(\sqrt{\frac{24k}{11m}} t\right)
		\right)}
	\tag{8}
	\label{eq:j17_p1_8}
\end{equation*}
Of course, using energies would produce the same result.
%	\newpage
\Solution{2}
We will begin with the second part of the question. The flipping is performed uniformly:
\begin{equation*}
	\boxed{\omega = \frac{\pi}{\tau}}
	\tag{$1$}
	\label{eq:j17_p2_1}
\end{equation*}

\begin{center}
	\begin{tikzpicture}
		% Ellipses
		\draw[line] (3, 4) ellipse (2cm and 0.25cm);
		\draw[line, rotate around={-30:(9, 4)}] (9, 4) ellipse (2cm and 0.25cm);
		\node[point] at (3, 4) {};
		\node[point] at (9, 4) {};
		
		% B-fields
		\draw[line, -{Stealth}] (3, 4) -- (3, 6);
		\draw[line, -{Stealth}] (9, 4) -- (9, 6);
		
		\node[above] at (3, 6) {$\vv{B}$};
		\node[above] at (9, 6) {$\vv{B}$};
		
		% Normals
		\draw[line, -{Stealth}] (3, 4) -- (3, 5);
		\draw[line, -{Stealth}, rotate around={-30:(9, 4)}] (9, 4) -- (9, 5);
		
		\node[right] at (3, 5) {$\hat{n}$};
		\node[right] at (9.5, 4.85) {$\hat{n}$};
		
		% Angle
		\coordinate (A) at (9, 6);
		\coordinate (B) at (9, 4);
		\coordinate (C) at (9.5, 4.85);
		
		\pic [draw, line, angle radius=0.5cm] {angle=C--B--A};
		
		\node[left] at (9.05, 4.61) {$\varphi$};
	\end{tikzpicture}
\end{center}

We can directly calculate the induced EMF in the ring as a function of time:
\begin{equation*}
	\varepsilon 
	= 
	-\frac{d\Phi_B}{dt}
	=
	- \frac{d}{dt}
	\left(
	\vv{B} \cdot \vv{S}
	\right)
	=
	-\frac{d}{dt}
	\left(
	\pi r^2 B \cos(\varphi)
	\right)
	=
	\pi r^2 B \sin(\varphi) \frac{d\varphi}{dt}
	\quad\Rightarrow\quad
	\boxed{
	\varepsilon(t)
	=
	\pi r^2 \omega B \sin(\omega t)
	}
	\tag{$2$}
	\label{eq:j17_p2_2}
\end{equation*}
Using equation (\ref{eq:j17_p2_1}) and Ohm's law:
\begin{equation*}
	i(t) = \frac{\varepsilon(t)}{R}
	\quad\Rightarrow\quad
	\boxed{i(t) = \frac{\pi^2 r^2 B}{R\tau} \sin\left(\frac{\pi t}{\tau}\right)}
	\tag{$3$}
	\label{eq:j17_p2_3}
\end{equation*}
The charge is then given by:
\begin{equation*}
	\Delta q
	=
	\int_0^\tau i(t)\,dt
	=
	\frac{1}{R}
	\int_0^\tau \varepsilon(t)\,dt
	=
	\frac{\pi r^2 B}{R}
	\int_0^\pi \sin\varphi\,d\varphi
	\quad\Rightarrow\quad
	\boxed{
		\Delta q 
		= 
		\frac{2\pi r^2 B}{R} 
		\approx 1{,}57\,\mathrm{\mu C}
		}
	\tag{$4$}
	\label{eq:j17_p2_4}
\end{equation*}
Note that this result does not depend on how the ring is flipped, provided it starts and ends in the same orientations. Consequently, the exact time dependence of the current  $i(t)$ is irrelevant for the total charge transferred. Only the initial and final magnetic fluxes matter:
\begin{equation*}
	\varepsilon = R \frac{dq}{dt} = - \frac{d\Phi_B}{dt}
	\quad\Rightarrow\quad
	\left|\Delta q \right|
	= 
	\frac{\left| \Delta \Phi_B \right|}{R} = 
	\frac{\pi r^2 B - (- \pi r^2 B)}{R}
	\quad\Rightarrow\quad
	\boxed{
		\left| \Delta q \right|
		= 
		\frac{2\pi r^2 B}{R} 
		\approx 1{,}57\,\mathrm{\mu C}
	}
\end{equation*}
%	\newpage
\Solution{3}
Notice that the speed of sound is proportional to the square root of the temperature:
\begin{equation*}
	\left|
	\begin{aligned}
		& c = \sqrt{
			\frac{\gamma p}{\rho}
		}
		\\
		& p = \frac{\rho}{\mu} RT
	\end{aligned}
	\,\,\,
	\right\}
	c = \sqrt{
	\frac{\gamma RT}{\mu}
	}
	\sim
	\sqrt{T}
	\tag{$*$}
	\label{eq:j17_p3_mock_1}
\end{equation*}
Using the sound increase given in the problem and differentiating:
\begin{equation*}
	dc = \frac{1}{2} \sqrt{ \frac{\gamma R}{\mu T} } dT
	\quad\Rightarrow\quad
	\frac{\Delta c}{c} \approx
	\frac{\Delta T}{2T}
	\quad\Rightarrow\quad
	\eta = \frac{\Delta c_0}{c} = \frac{\Delta T_0}{2 T}
	\quad\Rightarrow\quad
	\boxed{T = \frac{\Delta T_0}{2 \eta} = 250\,\mathrm{K}}
	\tag{$1$}
	\label{eq:j17_p3_1}
\end{equation*}
where $\Delta c_0$ is the increase in the speed of sound resulting from a temperature increase of  $\Delta T_0$.

Now assume that the temperature has increased by some value $\Delta T$. The wavelength has changed:
\begin{equation*}
	\left|
	\begin{aligned}
		& \lambda = c(T) / \nu
		\\
		& \lambda' = c(T+\Delta T) / \nu
	\end{aligned}
	\right.
\end{equation*}
However the distance between the two points of interest remains unchanged:
\begin{equation*}
	\lambda N = \lambda' N'
	\quad\Rightarrow\quad
	\frac{N}{N'} = \frac{c(T+\Delta T)}{c(T)}
	\approx
	\frac{c(T) + \Delta c}{c(T)}
	=
	1 + \frac{\Delta c}{c}
	\stackrel{(\refeq{eq:j17_p3_1})}{=}
	1 + \frac{\Delta T}{2T}
\end{equation*}
A temperature increase would mean that the speed of sound increases, hence the number of wavelengths between the two points will decrease. Therefore, $N' < N$. For antiphase, we require $N' = N - 1/2$:
\begin{equation*}
	\frac{N}{N - 1/2} = 1 + \frac{\Delta T}{2T}
	\quad\Rightarrow\quad
	\boxed{\Delta T = \frac{2T}{2N - 1} \approx 0{,}303\,\mathrm{K}}
	\quad\text{or}\quad
	\boxed{\Delta T = \frac{\Delta T_0}{\eta (2N - 1)} \approx 0{,}303\,\mathrm{K}}
	\tag{$2$}
	\label{eq:j17_p3_2}
\end{equation*}

\textbf{\underline{Remarks:}} Set $\alpha = \sqrt{\gamma R / \mu}$ in (\refeq{eq:j17_p3_mock_1}). Then Taylor-expand the equation for the speed of sound:
\begin{equation*}
	c(T) = \alpha \sqrt{T} 
	=
	\alpha
	\left(
	\sqrt{T_0}
	+
	\frac{T - T_0}{2 \sqrt{T_0}}
	-
	\frac{ (T - T_0)^2}{8 T_0^{3/2}}
	+
	\text{etc.}
	\right)
\end{equation*}
where we have expanded about a particular temperature $T_0$. Using a first-order approximation when temperature differences are small ($T \approx T_0$) yields:
\begin{equation*}
	c(T) 
	\approx
	\alpha \sqrt{T_0} 
	+
	\alpha \frac{T - T_0}{2 \sqrt{T_0}}
	\quad\Rightarrow\quad
	\boxed{
	c(T) \approx \alpha \sqrt{T_0}
	\left(
	1 + \frac{1}{2T_0}
	\left(T - T_0\right)
	\right)
	}
\end{equation*}
This is analogous to the familiar linear approximation for the temperature dependence of resistance:
\begin{equation*}
	\boxed{R(T) = R_0 (1 + \kappa (T - T_0))}
\end{equation*}
As long as the temperature difference remains small, we can define a ``temperature coefficient'':
\begin{equation*}
	\kappa_{eff} = \frac{1}{2 T_0}
\end{equation*}
We can also quickly check from the Taylor expansion that:
\begin{equation*}
	\frac{\Delta c}{c_0} = \frac{\Delta T}{2 T_0}
	\quad\text{or}\quad
	\eta = \kappa_{eff} \Delta T_0
\end{equation*}
where $c_0 = c(T_0)$ and we have evaluated $\Delta c / c_0$ at $\Delta T = \Delta T_0$. Therefore, our ``temperature coefficient'' is:
\begin{equation*}
	\boxed{\kappa_{eff} = \frac{\eta}{\Delta T_0}}
\end{equation*}
The idea is that a locally constant temperature coefficient could have been directly inferred from the problem statement as an approximation that could have been used in place of an explicit knowledge of the dependence $c = \sqrt{\gamma p / \rho}$. In fact, we do not need this formula at all to solve the problem!
%	\newpage
\Solution{4}
If $\vv{F}_1$ and $\vv{F}_2$ are parallel, then there exists some vector $\vv{u}$ which is perpendicular to both:
\begin{equation*}
	\vv{u} \cdot \vv{F}_1 = 0
	\quad\text{and}\quad
	\vv{u} \cdot \vv{F}_2 = 0
\end{equation*}
Then for the resultant force we show:
\begin{equation*}
	\vv{u} \cdot \vv{F}_{res}
	=
	\vv{u} \cdot 
	\left(
	\vv{F}_1 + \vv{F}_2
	\right)
	=
	\vv{u} \cdot \vv{F}_1
	+
	\vv{u} \cdot \vv{F}_2
	=
	0
	\quad\Rightarrow\quad
	\boxed{\vv{u} \cdot \vv{F}_{res} = 0}
	\tag{$1$}
	\label{eq:j17_p4_1}
\end{equation*}
Since $\vv{F}_{res}$ is also perpendicular to $\vv{u}$, then we have shown that it is parallel to both $\vv{F}_1$ and $\vv{F}_2$. Given this, the vector diagram for $\vv{F}_{res} = \vv{F}_1 + \vv{F}_2$ degenerates to a line segment. Therefore:
\begin{equation*}
	\boxed{F_{res} = F_1 + F_2
	}
	\tag{$2$}
	\label{eq:j17_p4_2}
\end{equation*}
\begin{minipage}[t][145pt]{\textwidth}
	\vspace{0pt}
	\begin{minipage}[t][143pt]{0.5\textwidth}
	\vspace{0pt}
	\centering
	\begin{tikzpicture}
		% Forces
		\node[point] at (3, 6) {};
		\node[point] at (6, 10) {};
		
		\node[below left] at (3, 6) {$A$};
		\node[above] at (6, 10) {$B$};
		
		\draw[line, -{Stealth}] (3, 6) -- (5.5, 6);
		\draw[line, -{Stealth}] (6, 10) -- (7.5, 10);
		
		\node[right] at (5.5, 6) {$\vv{F_1}$};
		\node[right] at (7.5, 10) {$\vv{F_2}$};
		
		% Resultant force
		\draw[line, -{Stealth}] (5.25, 9) -- (9, 9);
		\node[point] at (5.25, 9) {};
		\node[right] at (9, 9) {$\vv{F}_{res}$};
		
		% CM line
		\draw[line, dashed] (3, 7.5) -- (8, 7.5);
		
		\node[above] at (7, 7.5) {\text{CM line of}};
		\node[below] at (7, 7.5) {\text{action}};
		
		% Triangle
		\draw[line] (3, 6) -- (6, 10);
		\draw[line, dashed] (3, 6) -- (3, 10);
		\draw[line, dashed] (3, 10) -- (6, 10);
		\draw[line, dashed] (5.25, 9) -- (3, 9);
		
		\node[point] at (4.125, 7.5) {};
		\node[point] at (3, 7.5) {};
		\node[point] at (3, 10) {};
		\node[point] at (3, 9) {};
		
		\node[below] at (4.2, 7.5) {$C$};
		\node[left] at (3, 7.5) {$M$};
		\node[left] at (3, 10) {$N$};
		\node[above right] at (3, 9) {$T$};
		\node[above] at (5.2, 9) {$S$};
		
		% Brackets and dimensions
		\draw[bracket] (3, 6) -- (3, 7.5)
		node[midway, left=5pt] {$x$};
		\draw[bracket] (3, 7.5) -- (3, 10)
		node[midway, left=5pt] {$y$};
		\draw[line, Stealth-Stealth] (3.5, 7.55) -- (3.5, 8.95)
		node[midway, right] {$c$};
	\end{tikzpicture}
\end{minipage}
	\begin{minipage}[t][143pt]{0.49\textwidth}
		\vspace{2pt}
		For the last part consider the shown diagram. Assume that the distances from the lines of action of $\vv{F}_1$ and $\vv{F}_2$ to the line of action of the centre of mass are $x$ and $y$ respectively. In addition, let the distance from the line of action of $\vv{F}_{res}$ to the centre of mass line of action be $c$. From the torque law:
		\begin{equation*}
			F_2 y - F_1 x = F_{res} c = (F_1 + F_2) c
		\end{equation*}
		Some algebra yields:
		\begin{equation*}
			F_1 (x + c) = F_2 (y - c)
		\end{equation*}
	\end{minipage}
\end{minipage}
Notice that $x + c = AT$ and $y - c = TN$. Then we have:
\begin{equation*}
	\boxed{\frac{AT}{TN} = \frac{F_2}{F_1} = \frac{AS}{SB}}
	\tag{$3$}
	\label{eq:j17_p4_3}
\end{equation*}
where the last equality follows from Thales' Theorem.
%	\newpage
\Solution{5}
We treat each part separately:
\begin{description}[leftmargin=!, labelwidth=\widthof{$\bullet$}]
	\item[$\bullet$] \textbf{\underline{Electron Starting At Infinity:}} We will use conservation of energy as follows:
	\begin{center}
	\begin{tikzpicture}
		% Ellipse
		\draw[line] (10, 4) ellipse (0.25cm and 2cm);
		\node[point] at (10, 4) {};
		
		% Electron
		\node[point] at (2, 4) {};
		\node[above] at (2, 4) {$e^-$};
		\node[below] at (2, 4) {$\infty$};
		
		% Vector and Dashed Line
		\draw[line, dashed] (2, 4) -- (10, 4);
		\draw[line, -{Stealth}] (2, 4) -- (4, 4)
		node[above=1pt] {$\vv{v} = \vv{0}$};
	\end{tikzpicture}
\end{center}
	\begin{equation*}
		E_k^{\infty} + E_p^{\infty} = E_k + E_p
		\quad\Rightarrow\quad
		0 = \frac{1}{2} m_e v^2 - \frac{ke (\lambda 2 \pi R)}{R}
		\quad\Rightarrow\quad
		\frac{1}{2} m_e v^2 = \frac{e\lambda}{2\varepsilon_0}
		\quad\Rightarrow\quad
		\boxed{
			v = \sqrt{
				\frac{e\lambda}{m_e \varepsilon_0}
			}
		}
		\tag{$1$}
		\label{eq:j17_p5_1}
	\end{equation*}
	where we have used the following assumptions:
	\begin{description}[leftmargin=!, labelwidth=\widthof{$\rightarrow$}]
		\item[$\rightarrow$] Gravitational potential energy can be neglected as a whole;
		\item[$\rightarrow$] $E_p^{\infty} \rightarrow 0$ at the starting moment, since the electron is far away from the ring;
		\item[$\rightarrow$] $E_k^{ring} \rightarrow 0$ when the electron passes through the centre of the ring, since $M_{ring} >> m_e$;
	\end{description}
	\item[$\bullet$] \textbf{\underline{Electron Starting At Finite Distance}}
	\begin{center}
	\begin{tikzpicture}
		% Ellipse
		\draw[line] (10, 4) ellipse (0.25cm and 2cm);
		\node[point] at (10, 4) {};
		
		% Electron
		\node[point] at (2, 4) {};
		\node[above] at (2, 4) {$e^-$};
		
		% Vector and Dashed Line
		\draw[line, dashed] (2, 4) -- (10, 4);
		\draw[line, -{Stealth}] (2, 4) -- (4, 4)
		node[above=1pt] {$\vv{v} = \vv{0}$};
		\draw[bracket] (10, 3.95) -- (2, 3.95)
		node[midway, below=5pt] {$\eta R$};
	\end{tikzpicture}
\end{center}
	
	In this case only $E_k^{ring} \rightarrow 0$, so by conservation of energy:
	\begin{equation*}
		E_k^{0} + E_p^{0} = E_k + E_p
		\,\,\Rightarrow\,\,
		-\frac{k e (\lambda 2 \pi R)}{\sqrt{R^2 + \eta^2 R^2}} 
		=
		\frac{1}{2} m_e v^2 - \frac{ke(\lambda 2 \pi R)}{R}
		\,\,\Rightarrow\,\,
		\boxed{
			v
			=
			\sqrt{
				\frac{e \lambda}{m_e \varepsilon_0}
				\left(
				1 - \frac{1}{\sqrt{\eta^2 + 1}}
				\right)
			}
		}
		\tag{$2$}
		\label{eq:j17_p5_2}
	\end{equation*}
	We can see that taking the limit of (\ref{eq:j17_p5_2}) at $\eta \rightarrow +\infty$ restores (\ref{eq:j17_p5_1}):
	\begin{equation*}
		\lim\limits_{\eta \rightarrow + \infty}
		\left(
		\sqrt{
			\frac{e \lambda}{m_e \varepsilon_0}
			\left(
			1 - \frac{1}{\sqrt{\eta^2 + 1}}
			\right)
		}
		\right)
		=
		\sqrt{
			\frac{e \lambda}{m_e \varepsilon_0}
		}
	\end{equation*}
	\item[$\bullet$] \textbf{\underline{The Period Of Small Oscillations}}
	
	To find the period of small oscillations. consider the electron a distance $x$ away from the centre of the ring along its axis. The potential energy, using Bernoulli's approximation, of the system is then:
	\begin{equation*}
		E_p(x) = - \frac{
			ke (\lambda 2 \pi R)
		}{\sqrt{R^2 + x^2}}
		\approx
		\frac{\lambda e}{2 \varepsilon_0}
		\left(
		\frac{x^2}{2R^2} - 1
		\right)
		\quad\rightarrow\quad
		\boxed{k_{eff} = \frac{\lambda e}{2\varepsilon_0 R^2}}
		\tag{$*$}
		\label{eq:j17_p5_mock1}
	\end{equation*}
	where the constant term in the potential energy expression is just the potential energy of the system at equilibrium. As the ring is massive enough to neglect its motion, then the period of small oscillations is:
	\begin{equation*}
		T = 2 \pi \sqrt{\frac{m_e}{k_{eff}}}
		\quad\Rightarrow\quad
		\boxed{T = 2 \pi R \sqrt{\frac{2 m_e \varepsilon_0}{\lambda e}}}
		\tag{$3$}
		\label{eq:j17_p5_3}
	\end{equation*}
\end{description}

\newpage
\begin{description}[leftmargin=!, labelwidth=\widthof{$\bullet$}]
	\item[$\bullet$] \textbf{\underline{Force Approximation}}
	
	We require the magnitude of the difference to be no more than 1\% w.r.t. the linearized force:
	\begin{equation*}
		\frac{|F_x - F_{lin}|}{F_{lin}} \leq 0{,}01
		\quad\Rightarrow\quad
		\left|
		1 - \frac{F_x}{F_{lin}}
		\right|
		\leq
		0{,}01
	\end{equation*}
	
	\begin{minipage}[t][117pt]{\textwidth}
		\vspace{0pt}
		\begin{minipage}[t][116pt]{0.35\textwidth}
	\vspace{0pt}
	\centering
	\begin{tikzpicture}
		% Ellipse
		\draw[line] (10, 4) ellipse (0.25cm and 2cm);
		\node[point] at (10, 4) {};
		\draw[line, dashed] (10, 4) -- (11, 4);
		\draw[line, dashed] (10, 6) -- (11, 6);
		\draw[line, Stealth-Stealth] (11, 4.05) -- (11, 5.95)
		node[midway, right] {$R$};
		
		% Electron
		\node[point] at (7, 4) {};
		\node[left] at (7, 4) {$e^-$};
		\draw[line, dashed] (7, 4) -- (10, 4)
		node[midway, below] {$x$};
		\draw[line] (7, 4) -- (10, 6)
		node[midway, above] {$r$};
		
		% Force
		\draw[line, -{Stealth}] (7, 4) -- (8, 4)
		node[below] {$F_x$};
	\end{tikzpicture}
\end{minipage}
		\begin{minipage}[t][116pt]{0.64\textwidth}
			\vspace{0pt}
			We can easily get expressions for the forces from the diagram:
			\begin{equation*}
				\left|
				\begin{aligned}
					& F_x = \frac{ke (\lambda 2 \pi R)}{r^2} \cdot \frac{x}{r}
					=
					\frac{\lambda e}{2 \varepsilon_0 R^2}
					\cdot
					\frac{x}{(1 + (x/R)^2)^{3/2}}
					\\
					& F_{lin} = k_{eff} x = \frac{\lambda e x}{2\varepsilon_0 R^2}
				\end{aligned}
				\right.
			\end{equation*}
			We can now plug these in the above expression:
			\begin{equation*}
				1 - \frac{1}{(1 + (x/R)^2)^{3/2}} \leq 0{,}01
			\end{equation*}
		\end{minipage}
	\end{minipage}
	
	\vspace{1em}
	where the modulus sign has now been dropped. Finishing the calculation:
	\begin{equation*}
		\boxed{x \leq R \sqrt{0{,}99^{-2/3} - 1}}
		\quad\text{or}\quad
		\boxed{x \leq 0{,}082 R}
		\tag{$4$}
		\label{eq:j17_p5_4}
	\end{equation*}
	The force approximation becomes $1\%$ inaccurate at about $8{,}2\%$ of the ring radius.
	\item[$\bullet$] \textbf{\underline{Period Approximation}}
	
	Let the amplitude be $A = \gamma R$ and release the electron from rest from it. By conservation of energy:
	\begin{equation*}
		- \frac{k e (\lambda 2 \pi R)}{\sqrt{A^2 + R^2}} 
		=
		\frac{1}{2} m v^2 
		-
		\frac{k e (\lambda 2 \pi R)}{\sqrt{A^2 + x^2}}
		\quad\Rightarrow\quad
		\boxed{
		(v(x))^2 = 
			\frac{e \lambda}{m_e \varepsilon_0}
			\left(
			\frac{1}{\sqrt{1 + (x/R)^2}}
			-
			\frac{1}{\sqrt{1 + \gamma^2}}
			\right)
		}
		\tag{$5a$}
		\label{eq:j17_p5_5a}
	\end{equation*}
	Now use the Taylor expansion in (\ref{eq:j17_p5_5a}) for the following (taken from the Appendix):
	\begin{equation*}
		\frac{1}{\sqrt{1+x}} = 1 - \frac{1}{2} x + \frac{3}{8} x^2 + \mathcal{O}(x^3)
		\,\Rightarrow\,
		(v(x))^2 = \frac{e\lambda}{m_e \varepsilon_0}
		\left(
		1 -
		\frac{1}{\sqrt{1 + \gamma^2}}
		-
		\frac{1}{2} \left(\frac{x}{R}\right)^2
		+
		\frac{3}{8} \left(\frac{x}{R}\right)^4
		+
		\mathcal{O}\left(\left(\frac{x}{R}\right)^6\right)
		\right)
	\end{equation*}
	Introduce the constant $K$ as follows and rewrite:
	\begin{equation*}
		\left|
		\begin{aligned}
			& K = 1 - \frac{1}{\sqrt{1 + \gamma^2}} 
			=
			\frac{1}{2} \gamma^2 
			-
			\frac{3}{8} \gamma^4
			+
			\mathcal{O}(\gamma^6)
			\rowstrut
			\\
			& v(x) =
			\sqrt{
			\frac{e\lambda}{m_e \varepsilon_0}
			\left(
			K - \frac{1}{2} \left(\frac{x}{R}\right)^2
			+
			\frac{3}{8} \left(\frac{x}{R}\right)^4
			+
			\mathcal{O}\left(\left(\frac{x}{R}\right)^6\right)
			\right)
			}
			\rowstrut
		\end{aligned}
		\right.
		\tag{$5b$}
		\label{eq:j17_p5_5b}
	\end{equation*}
	Finding the period requires evaluating the following integral\footnote{Note that $dt = - dx / v(x)$ as $dt > 0$ and $dx < 0$.}:
	\begin{equation*}
		T = 4\int\limits_{0}^{T/4} dt =  
		- 4 \int\limits_{A}^{0} \frac{dx}{v(x)}
		=
		4 \sqrt{\frac{m_e \varepsilon_0}{e \lambda}}
		\int\limits_{0}^{A}
		\left(
		K - \frac{1}{2} \left(\frac{x}{R}\right)^2
		+
		\frac{3}{8} \left(\frac{x}{R}\right)^4
		+
		\mathcal{O}\left(\left(\frac{x}{R}\right)^6\right)
		\right)^{-1/2}
		dx
	\end{equation*}
	We will now expand $K$ first in the integrand and later expand the integrand itself (see remark):
	\begin{equation*}
		T = 
		4 \sqrt{\frac{m_e \varepsilon_0}{e \lambda}}
		\int\limits_{0}^{A}
		\left(
		\frac{1}{2} \gamma^2 - \frac{3}{8} \gamma^4 - \frac{1}{2} \left(\frac{x}{R}\right)^2
		+
		\frac{3}{8} \left(\frac{x}{R}\right)^4
		+
		\mathcal{O}(\gamma^6)
		\right)^{-1/2}
		dx
		\tag{$**$}
		\label{eq:j17_p5_mock2}
	\end{equation*}
	Here we used the fact that $\gamma = A/R$, i.e. it follows that $x / R = \mathcal{O}(\gamma)$ over the interval of integration, so we simply included the $\mathcal{O}((x/R)^6)$ term inside $\mathcal{O}(\gamma^6)$. Factor out a zero at $x = A$ (also see remark):
	\begin{equation*}
		T
		=
		4 \sqrt{\frac{m_e \varepsilon_0}{e \lambda}}
		\int\limits_{0}^{A}
		\left(
		\frac{A^2 - x^2}{2R^2}
		-
		\frac{3(A^4 - x^4)}{8R^4}
		+
		\mathcal{O}(\gamma^6)
		\right)^{-1/2} dx
	\end{equation*}
	\begin{equation*}
		T
		=
		4 \sqrt{\frac{m_e \varepsilon_0}{e \lambda}}
		\int\limits_{0}^{A}
		\frac{R \sqrt{2}}{\sqrt{A^2 - x^2}}
		\left(
		1
		-
		\frac{3(A^2 + x^2)}{4R^2}
		+
		\mathcal{O}(\gamma^4)
		\right)^{-1/2} dx
	\end{equation*}
	Finally, we can approximate the expression in the brackets:
	\begin{equation*}
		T
		\approx
		4 \sqrt{\frac{m_e \varepsilon_0}{e \lambda}}
		\int\limits_{0}^{A}
		\frac{R \sqrt{2}}{\sqrt{A^2 - x^2}}
		\left(
		1
		+
		\frac{3(A^2 + x^2)}{8R^2}
		\right) 
		dx
		=
		4R \sqrt{\frac{2m_e \varepsilon_0}{e \lambda}}
		\int\limits_{0}^{A}
		\left(
		\frac{1}{\sqrt{A^2 - x^2}}
		+
		\frac{3 (A^2 + x^2)}{8R^2 \sqrt{A^2 - x^2}}
		\right)
		dx
	\end{equation*}
	Evaluating these integrals yields:
	\begin{equation*}
		\left|
		\begin{aligned}
			& \int\limits_{0}^{A} \frac{dx}{\sqrt{A^2 - x^2}} = \frac{\pi}{2}
			\rowstrut
			\\
			& \int\limits_{0}^{A} \frac{A^2 + x^2}{\sqrt{A^2 - x^2}} dx
			=
			\frac{3 \pi A^2}{4}
		\end{aligned}
		\,\,
		\right\}
		\quad
		\boxed{
		T
		\approx
		2 \pi R \sqrt{\frac{2 m_e \varepsilon_0}{e \lambda}}
		\left(1 + \frac{9A^2}{16R^2}\right)
		}
		\tag{$6$}
		\label{eq:j17_p5_6}
	\end{equation*}
	Notice that this is larger than the period for small oscillations from equation (\ref{eq:j17_p5_3}). Therefore we require:
	\begin{equation*}
		\frac{
		|T - T_0|
		}{T_0} \leq 0{,}01
		\quad\Rightarrow\quad
		\left|
		1 - \frac{T}{T_0}
		\right| \leq 0{,}01
		\quad\Rightarrow\quad
		\frac{9A^2}{16R^2} \leq 0{,}01
	\end{equation*}
	Solving for $A$ gives:
	\begin{equation*}
		\boxed{A \leq \frac{4}{30} R}
		\quad\Rightarrow\quad
		\boxed{A \leq 0{,}133 R}
		\tag{$7$}
		\label{eq:j17_p5_7}
	\end{equation*}
	The period approximation becomes $1\%$ inaccurate at about $13{,}3\%$ of the ring radius.
\end{description}
\textbf{\underline{Remarks:}} Getting to expand an expression under the integral sign can get tricky. Take for example the expression from (\ref{eq:j17_p5_mock2}) and factor out the first term, then apply the Taylor expansion:
	\begin{equation*}
		\begin{aligned}
			T 
			&= 
			4 \sqrt{\frac{m_e \varepsilon_0}{e \lambda}}
			\int\limits_{0}^{A}
			\left(
			\frac{1}{2} \gamma^2 - \frac{3}{8} \gamma^4 - \frac{1}{2} \left(\frac{x}{R}\right)^2
			+
			\frac{3}{8} \left(\frac{x}{R}\right)^4
			+
			\mathcal{O}(\gamma^6)
			\right)^{-1/2}
			dx
			\approx
			\rowstrut
			\\
			&\approx
			4 \sqrt{\frac{m_e \varepsilon_0}{e \lambda}}
			\int\limits_{0}^{A}
			\frac{\sqrt{2}}{\gamma}
			\left(
			1 - \frac{3}{4} \gamma^2 
			-
			\frac{1}{\gamma^2} \left(\frac{x}{R}\right)^2 
			+
			\frac{3}{4\gamma^2} \left(\frac{x}{R}\right)^4
			\right)^{-1/2} dx
			\approx
			\rowstrut
			\\
			&\approx
			4 \sqrt{\frac{m_e \varepsilon_0}{e \lambda}}
			\int\limits_{0}^{A}
			\frac{R \sqrt{2}}{A}
			\left(
			1 + \frac{3A^2}{8R^2} + \frac{x^2}{2A^2}
			-
			\frac{3x^4}{8A^2 R^2}
			\right)
			dx =
			\rowstrut
			\\
			&= 
			\sqrt{\frac{m_e \varepsilon_0}{e \lambda}}
			\left(
			\frac{14 \sqrt{2}}{3} R + \frac{6 \sqrt{2} A^2}{5R}
			\right)
		\end{aligned}
\end{equation*}
which is a radically different expression from what we got. So where did we go wrong?

Notice that the original expression from (\ref{eq:j17_p5_mock2}) is $1 / v(x)$ and it admits $v(A) = 0$ at the end of the integration interval. In this example, when we factored out $\gamma^2 / 2$, we lost this property. We get on the second line of the above calculation an expression of the form $(1 + u(x, \gamma))^{-1/2}$ the binomial expansion of which requires $|x| < 1$. However, $u(A, \gamma) = -1$ (because $v(A) = 0$) which violates this condition whenever $x \rightarrow A$. 

Expanding in this way destroys the zero of \(v(x)\) at the turning point and therefore destroys the corresponding integrable singularity of \(1/v(x)\). Although the expansion may be valid away from the endpoint, it is not uniform on the entire interval \(0\le x\le A\), and hence it cannot in general be integrated term by term to obtain the correct asymptotic expansion of the period.

This is why, in the previous derivation, it was important to first factor out \(A^2-x^2\). Doing so preserves the turning-point behaviour explicitly:

\begin{equation*}
	\frac{1}{v(x)}
	\sim
	\frac{1}{\sqrt{A^2-x^2}}
	\left(1+\mathcal{O}\!\left(\frac{A^2}{R^2}\right)\right)
\end{equation*}

after which the remaining expansion is uniformly small over the integration interval.

%	\newpage
\Solution{6}
We have two conditions present:
\begin{description}
	\item[$\rightarrow$] The vessels are thermally insulated, so $Q = 0$;
	\item[$\rightarrow$] The gas does no work on the vacuum, so $A = 0$;
\end{description}
By the First Law of Thermodynamics:
\begin{equation*}
	Q = \Delta U + A
	\quad\Rightarrow\quad
	\Delta U = 0
	\quad\Rightarrow\quad
	\boxed{\Delta T = 0}
	\tag{$1$}
	\label{eq:j17_p6:1}
\end{equation*}
where the latter follows from the fact that for an ideal gas $U = U(T)$.

Although the actual free expansion is irreversible, entropy is a function of state, so we can calculate $\Delta S$ using a convenient reversible process or combination of processes. The easiest choice to make for this case, due to equation (\ref{eq:j17_p6:1}), is a \emph{\underline{reversible isothermal process}}\footnote{It is instructive to choose a different process and verify that the same result is achieved. For example, the answer for $\Delta S$ can be verified using an isobaric process and an isochoric process, ``connecting'' the initial and final states.}:
\begin{equation*}
	\Delta S = \int\limits_1^2 \frac{dQ_{rev}}{T}
	=
	\int\limits_{V_1}^{V_2}
	\nu R \frac{dV}{V}
	=
	\nu R \ln\left(\frac{V_2}{V_1}\right)
	\quad\Rightarrow\quad
	\boxed{\Delta S = \nu R \ln(n) \approx 20{,}1\,\mathrm{J/K}}
	\tag{$2$}
	\label{eq:j17_p6_2}
\end{equation*}
Note that this does not mean the original process is not isothermal: it is an \emph{\underline{irreversible isothermal process}}.

\textbf{\underline{Remarks:}} Here entropy increases, despite $Q = 0$. It is important to understand that:
\begin{equation*}
	dS = \frac{dQ_{rev}}{T}
\end{equation*}
holds only for a \emph{\underline{reversible}} process. More generally, for an arbitrary process:
\begin{equation*}
	\boxed{dS \geq \frac{dQ}{T}}
\end{equation*}
with equality only in the reversible case.

In the present problem, the entropy increase does not originate from heat supplied to the gas. Instead, it is a consequence of the \emph{\underline{irreversibility}} of the free expansion: the gas initially occupies only one of the vessels, while after opening the valve it spontaneously occupies the larger available volume. The number of accessible microscopic states of the gas consequently increases. Therefore, for one particle we can state:
\begin{equation*}
	\Omega \sim V
\end{equation*}
where the momentum-space contribution is unchanged, since  $\Delta T = 0$.

For $N$ particles, the number of accessible microstates is proportional to:
\begin{equation*}
	\Omega \sim V^N
	\quad\Rightarrow\quad
	\frac{\Omega_2}{\Omega_1} = \left(\frac{V_2}{V_1}\right)^N
	=
	n^N
\end{equation*}
Finally, using Boltzmann's Entropy Relation:
\begin{equation*}
	S = k \ln(\Omega)
	\quad\Rightarrow\quad
	\Delta S = k \ln\left(\frac{\Omega_2}{\Omega_1}\right)
	=
	Nk \ln(n)
	\quad\Rightarrow\quad
	\boxed{\Delta S = \nu R \ln(n) \approx 20{,}1\,\mathrm{J/K}}
\end{equation*}

This also illustrates why, when calculating the entropy change of an irreversible process, one may replace the actual process by any convenient reversible process connecting the same initial and final equilibrium states. The reversible process is used only as a means of calculating the state-function difference $\Delta S$; it is not a description of what actually happens to the gas.
	
	% Национална Олимпиада По Физика 2025 - Тема 10 Клас - Задача 3, Част 2
	% Национална Олимпиада По Физика 2013 - Тема 10-12 Клас - Задача 2, Част 1
	% Национална Олимпиада По Физика 2013 - Тема 10-12 Клас - Задача 2, Част 2
	% 3800 Задач Физике - 4.39
	% 3800 Задач Физике - 11.231 + 11.232 + 11.233
	% Irodov 2.148
\end{document}